% Author: Izaak Neutelings (March 2024)
% Description: Lead-lead collisions
\documentclass[border=3pt,tikz]{standalone}
\usepackage{amsmath}
\usepackage{listofitems} % to create arrays with \readlist
\usepackage{ifthen} % for \whiledo
%\usetikzlibrary{arrows.meta} % for arrow size

% COLORS
\colorlet{myred}{red!75!black}
\colorlet{myblue}{blue!70!black}

% STYLES
\tikzset{
  >=latex, % for LaTeX arrow head
  ball/.style={ultra thin,draw=black,fill=\getcol{#1},
               circle,inner sep=1pt,minimum size=15pt,scale=0.8},
  ball/.default=1,
}
\def\getcol#1{\ifodd#1 blue!20 \else red!20 \fi}
\def\drawgrid#1{
  \pgfmathsetmacro\xmax{ceil(#1)+0.6}
  \pgfmathsetmacro\xmin{0.3-\xmax}
  \draw[black!20] (\xmin,\xmin) grid (\xmax,\xmax);
  \draw[thick,->] (\xmin,0) -- (\xmax,0);
  \draw[thick,->] (0,\xmin) -- (0,\xmax);
  \draw[black!20] (0,0) circle(#1);
}

\begin{document}


% NUCLEUS - random radius & angle
\def\R{2.0}
\def\Nlay{5} % number of rings/layers
\def\N{100}
\begin{tikzpicture}
  \drawgrid{\R}
  
  % NUCLEUS
  \foreach \i [evaluate={
    \r=\R*sqrt((rand+1)/2); % ball radius
    \a=180*rand % ball angle
  }] in {1,...,\N}{
    \node[ball=\i] at (\a:\r) {\i};
  }
  
\end{tikzpicture}


% NUCLEUS - uniform rings
\begin{tikzpicture}
  \message{^^JNucleus with uniform rings}
  \drawgrid{\R}
  
  % NUCLEUS
  \foreach \lay [evaluate={
    \r=\R*\lay/\Nlay; % ring radius
    \Nring=int(round(2*\lay*\N/(\Nlay*(\Nlay+1)))) % number of balls per ring
  }] in {\Nlay,...,1}{
    \message{^^J N=\N, Nlay=\Nlay, Nring=\Nring, r=\r}
    \foreach \i [evaluate={
      \a=360*(\i/\Nring); % ball angle
    }] in {1,...,\Nring}{
      \node[ball=\i] at (\a:\r) {\i};
    }
  }
  
\end{tikzpicture}


% NUCLEUS - uniform rings + random offset
\begin{tikzpicture}
  \message{^^JNucleus with uniform rings + random offset}
  %\def\R{1.8}
  \drawgrid{\R}
  
  % NUCLEUS
  \foreach \lay [evaluate={
    \r=\R*(\lay-0.1)/\Nlay; % ring radius
    \aR=3*rand; % angular offset
    \Nring=int(round(2*\lay*\N/(\Nlay*(\Nlay+1)))) % number of balls per ring
  }] in {\Nlay,...,1}{
    \message{^^J N=\N, Nlay=\Nlay, Nring=\Nring, r=\r}
    \foreach \i [evaluate={
      \aB=\aR+360*(\i/\Nring)+2*rand/\Nring; % ball angle plus random offset
      \rB=\r+0.04*rand % ball radius plus random radial offset
    }] in {1,...,\Nring}{
      \node[ball=\i] at (\aB:\rB) {\i};
      %\fill[red] ({\aR+360*(\i/\Nring)}:\r) circle(0.5pt);
    }
  }
  
\end{tikzpicture}


% NUCLEUS - shuffle balls
\def\shufflelist#1{
  %\message{^^JShuffle list}
  \pgfmathsetmacro\tmplist{random(1,#1)} % first random item
  \readlist*\shuffledlist\tmplist % convert to array
  \newboolean{unique}
  \whiledo{\shuffledlistlen<#1}{ % try until list is of length \N=#1
    \pgfmathsetmacro\rx{random(1,#1)} % generate new random item
    \setboolean{unique}{true}
    \foreachitem \x \in \shuffledlist{ % loop over previous elements
      \ifnum \x=\rx % look if new number in list
        \setboolean{unique}{false} % already in list
        \breakforeach
      \fi
    }
    \ifthenelse{\boolean{unique}}{ % found unique item !
      \xdef\tmplist{\tmplist,\rx} % add new unique item
      \readlist*\shuffledlist\tmplist % convert to array
    }{}
  }
}
\begin{tikzpicture}[xscale=1]
  \message{^^JNucleus with shuffling balls}
  %\def\R{1.8}
  \drawgrid{\R}
  \pgfmathsetseed{1234567890}
  
  % NUCLEUS
  \foreach \lay [evaluate={
    \r=\R*(\lay-0.1)/\Nlay; % ring radius
    \aR=3*rand; % angular offset
    \Nring=int(round(2*\lay*\N/(\Nlay*(\Nlay+1)))) % number of balls per ring
  }] in {\Nlay,...,1}{
    \message{^^J N=\N, Nlay=\Nlay, Nring=\Nring, r=\r}
    \shufflelist{\Nring} % get shuffled list of sequence 1,...,\Nring
    \foreach \i [evaluate={
      \is=int(\shuffledlist[\i]); % shuffled index
      \aB=\aR+360*(\is/\Nring)+2*rand/\Nring; % ball angle plus random offset
      \rB=\r+0.04*rand % ball radius plus random radial offset
    }] in {1,...,\Nring}{
      \message{^^J is=\is, i=\i, a=\aB}
      \node[ball=\i] at (\aB:\rB) {\i};
      %\fill[red] ({\aR+360*(\i/\Nring)}:\r) circle(0.5pt);
    }
  }
  \node[ball=\i] at (0,0) {0}; % cover hole
  
\end{tikzpicture}


% NUCLEUS - 3D effect
\begin{tikzpicture}[xscale=1]
  \message{^^JNucleus with 3D effect}
  %\def\R{1.8}
  \drawgrid{\R}
  \pgfmathsetseed{1234567890}
  
  % NUCLEUS
  \pgfmathsetmacro\C{min(1/(4*\R),0.11)} % put balls at edges closer in to create 3D ball effect
  \pgfmathsetmacro\D{1-4*\C*\R} % discriminant
  \message{^^J R=\R, C=\C, D=1-4cR=\D (must be >0)}
  \pgfmathsetmacro\Rs{2*\R/(1+sqrt(\D)} % scaled total radius
  \foreach \lay [evaluate={
    \r=\Rs*(\lay-0.1)/\Nlay; % ring radius
    \rC=\r-\C*\r*\r; % reduce outer radii to create 3D effect
    \aR=3*rand; % angular offset
    \Nring=int(round(2*\rC*\N/(\Rs*(\Nlay+1)*(1-\C*\Rs*(2*\Nlay+1)/(3*\Nlay))))) % number of balls on this ring
  }] in {\Nlay,...,1}{
    \message{^^J N=\N, Nlay=\Nlay, Nring=\Nring, r=\r}
    \shufflelist{\Nring} % get shuffled list of sequence 1,...,\Nring
    \foreach \i [evaluate={
      \is=int(\shuffledlist[\i]); % shuffled index
      \aB=\aR+360*(\is/\Nring)+2*rand/\Nring; % ball angle plus random offset
      \rB=\rC+0.05*rand % ball radius plus random radial offset
    }] in {1,...,\Nring}{
      \message{^^J is=\is, i=\i, a=\aB}
      \node[ball=\i] at (\aB:\rB) {\i};
      %\fill[red] ({\aR+360*(\i/\Nring)}:\r) circle(0.5pt);
    }
  }
  \node[ball=\i] at (0,0) {0}; % cover hole
  
\end{tikzpicture}


\end{document}